%-----------------------------------
\section{O Trabalho de Conclusão de Curso em Sistemas de Informação}
\label{sec_tcc_si}

Sistemas de Informação é uma área que combina conhecimentos de computação e de administração para conceber, desenvolver e gerenciar soluções que apoiam pessoas e organizações. Por isso, os TCCs do curso podem ter caráter predominantemente técnico, como a construção de um sistema; empírico, como a avaliação de uma tecnologia com usuários; ou organizacional, como a análise de processos de uma empresa.

Independentemente do tipo, todo TCC deve apresentar um problema bem delimitado, objetivos verificáveis, um método explícito e resultados discutidos à luz da literatura \cite{wazlawick2014metodologia}. O que muda de um tipo para outro é o \emph{como}: os procedimentos adotados, os capítulos necessários e os elementos usados para apresentar os resultados.

A monografia segue a estrutura definida pela ABNT NBR 14724:2011, composta por elementos pré-textuais, textuais e pós-textuais. Os elementos textuais deste modelo, comuns a todos os tipos de trabalho, são apresentados no Quadro~\ref{qua_estrutura_comum}.

\begin{quadro}[htb]
	\centering
	\caption{Elementos textuais comuns a todos os tipos de TCC.}
	\label{qua_estrutura_comum}
	\small
	\begin{tabularx}{\textwidth}{|l|X|}
		\hline
		\textbf{Capítulo} & \textbf{Conteúdo esperado} \\ \hline
		Introdução & Contextualização, problema de pesquisa, objetivos, justificativa e organização do trabalho. \\ \hline
		Fundamentação teórica & Conceitos e tecnologias que sustentam o trabalho e trabalhos relacionados. \\ \hline
		Metodologia & Classificação da pesquisa, etapas, coleta e análise dos dados e aspectos éticos. \\ \hline
		Resultados e discussão & Apresentação dos resultados e sua interpretação frente aos objetivos e à literatura. \\ \hline
		Considerações finais & Atendimento aos objetivos, contribuições, limitações e trabalhos futuros. \\ \hline
	\end{tabularx}
	\fonte{Elaborado pelo autor.}
\end{quadro}

Alguns tipos de trabalho acrescentam um capítulo entre a metodologia e os resultados, como o capítulo de desenvolvimento do sistema (Seção~\ref{sec_tipo_desenvolvimento}) ou o de caracterização do caso (Seção~\ref{sec_tipo_estudo_caso}). Consulte também o regulamento de TCC do curso, que pode estabelecer requisitos adicionais.
